\section{Accent Robustness in Speech Processing}
\label{sec:accent-asr}

% TODO(Member 2): Owner: Member 2. Target approximately 1.7 pages.
% Purpose: synthesize evidence on regional/accented English ASR, including
% seen versus unseen accents, underrepresented groups, generalization,
% adaptation, and limitations.
% Evidence needed: approximately five personally read peer-reviewed papers.
% Comparisons needed: accent definitions, speaker-disjoint evaluation,
% training/test overlap, dataset balance, WER patterns, adaptation trade-offs,
% and which populations are missing or aggregated.
% Boundary: this section may establish accent sensitivity in ASR but must not
% present ASR WER as direct evidence of final translation quality.
% Transition in: explain why architecture evidence must be paired with speech
% variation evidence.
% Transition out: make explicit which ASR findings motivate, but do not answer,
% the speech-translation question in Section IV.
\collabplaceholder{Member 2}{Develop the accent-ASR synthesis using verified
sources. Compare datasets, evaluation conditions, generalization, adaptation,
strengths, and limitations, while preserving the boundary between recognition
and final translation evidence.}
